\documentclass[12pt,a4paper]{article}

\usepackage[utf8]{inputenc}
\usepackage[T1]{fontenc}
\usepackage{lmodern}
\usepackage{amsmath,amssymb,amsthm}
\usepackage[top=2.5cm,bottom=2.5cm,left=2.8cm,right=2.8cm]{geometry}
\usepackage{parskip}
\usepackage{booktabs}
\usepackage{xcolor}
\usepackage{hyperref}
\hypersetup{colorlinks=true, linkcolor=blue!60!black}

\newtheorem{theorem}{Théorème}[section]
\newtheorem{lemma}[theorem]{Lemme}
\newtheorem{corollary}[theorem]{Corollaire}
\newtheorem{definition}[theorem]{Définition}
\theoremstyle{remark}
\newtheorem{note}{Note}[section]

\newcommand{\Inc}{\ensuremath{\mathrm{Include}}}
\newcommand{\Exc}{\ensuremath{\mathrm{Exclude}}}
\newcommand{\vbar}{\bar{v}}

\begin{document}

\begin{center}
  {\LARGE\bfseries Théorie consolidée de la règle à 4 cas}\\[6pt]
  {\large Automate isolé, sélection de littéral, extension multi-bit --- base pour la Section~4 (AAAI)}
\end{center}
\vspace{0.5em}\hrule\vspace{1.5em}

\tableofcontents
\newpage

% ======================================================
\section{Cadre général}
% ======================================================

Soit $(X,Y)\in\{0,1\}^2$ un couple aléatoire de loi jointe $P(X,Y)$ arbitraire (pas
nécessairement IDENTITY/NOT), et $(X_t,Y_t)_{t\ge1}\overset{\text{i.i.d.}}{\sim}P$
la suite des observations présentées à la règle -- hypothèse maintenue dans
toute la suite du document, nécessaire pour que chaque automate soit une
chaîne de Markov homogène à probabilités de transition constantes $p^+,p^-$
(sans elle, si les observations sont corrélées dans le temps ou si $P$ dérive,
l'automate seul n'est pas garanti être une chaîne de Markov homogène). Deux
automates gouvernent le littéral $x$ et sa négation : $v$ (état
$s_1\in\{1,\ldots,2N\}$) et $\vbar$ (état $s_2$), avec
$\mathrm{included}(s)=\mathbf 1[s>N]$. La règle à 4 cas, appliquée à chaque
observation $(x,y)$ :

\begin{center}
\begin{tabular}{ccll}
\toprule
$x$ & $y$ & Action sur $v$ & Action sur $\vbar$ \\
\midrule
0 & 0 & $\texttt{versInclude}(S)$ & $\texttt{versExclude}(S)$ \\
0 & 1 & $\texttt{versExclude}(a)$ & rien \\
1 & 0 & $\texttt{versExclude}(S)$ & $\texttt{versInclude}(S)$ \\
1 & 1 & rien & $\texttt{versExclude}(a)$ \\
\bottomrule
\end{tabular}
\end{center}

avec $S,a\in(0,1]$. En agrégeant sur $(X,Y)$, les probabilités de transition
intérieures de $v$ et $\vbar$ s'écrivent, pour \emph{toute} loi jointe $P$ :
\begin{align}
  p^+ &= S\cdot P(X{=}0,Y{=}0), &
  p^- &= S\cdot P(X{=}1,Y{=}0) + a\cdot P(X{=}0,Y{=}1); \label{eq:pv}\\
  \bar p^+ &= S\cdot P(X{=}1,Y{=}0), &
  \bar p^- &= S\cdot P(X{=}0,Y{=}0) + a\cdot P(X{=}1,Y{=}1). \label{eq:pvbar}
\end{align}

\begin{note}[Lien avec la paramétrisation $(c,p_1,p_0)$]
Le cadre IDENTITY/NOT déjà utilisé dans \textit{Démonstration de convergence}
($c{=}P(X{=}1)$, cette référence notant $a{=}P(Y{=}1\mid X{=}1)$ et
$b{=}P(Y{=}1\mid X{=}0)$ -- renommés ici $p_1,p_0$ pour éviter toute collision
avec le paramètre de règle $a$ introduit ci-dessous, sans lien avec ces
probabilités de bruit) est le cas particulier
$P(0,0){=}(1{-}c)(1{-}p_0)$, $P(1,0){=}c(1{-}p_1)$, $P(0,1){=}(1{-}c)p_0$,
$P(1,1){=}c\,p_1$. Toutes les formules ci-dessous se réduisent aux formules
déjà publiées dans \textit{Démonstration de convergence} en substituant ces
expressions --- vérifié terme à terme, aucune contradiction.
\end{note}

% ======================================================
\section{Théorème 1 --- loi stationnaire exacte (cas général)}
% ======================================================

\begin{theorem}[Loi stationnaire, automate isolé]\label{thm:stationnaire}
Soit $\rho = p^+/p^-$ (resp.\ $\bar\rho=\bar p^+/\bar p^-$), avec $p^+,p^->0$
(cas ergodique, aucune des deux directions n'est bloquée). La chaîne $(v(t))_t$
est irréductible et apériodique sur $\{1,\ldots,2N\}$, et sa loi stationnaire
est géométrique, de normalisation exacte :
\[
  \pi(k) = \frac{(1-\rho)\,\rho^{k-1}}{1-\rho^{2N}} \ (\rho\ne1), \qquad
  \pi(k) = \frac1{2N} \ (\rho=1), \qquad
  \boxed{\pi(\Inc) = \sum_{k>N}\pi(k) = \frac{\rho^N}{1+\rho^N}.}
\]
Si $\rho>1$, $\pi(\Inc)>\tfrac12$ et $\pi(\Inc)\xrightarrow[N\to\infty]{}1$, avec
$1-\pi(\Inc)\le\rho^{-N}$ ; si $\rho<1$, $\pi(\Exc)\to1$ symétriquement. De plus,
par le théorème ergodique pour les chaînes de Markov irréductibles apériodiques
à espace d'états fini,
\[
  \frac1T\sum_{t=1}^{T}\mathbf 1[v_t>N] \xrightarrow[T\to\infty]{\text{p.s.}} \pi(\Inc),
\]
ce qui relie directement la loi stationnaire à la métrique expérimentale
« fraction du temps passé dans la bonne action ». Idem pour $\vbar$ avec $\bar\rho$.
\end{theorem}

\begin{proof}
Irréductibilité : chaîne de naissance-mort à transitions intérieures
strictement positives ($p^+,p^->0$). Apériodicité : \emph{ne dépend pas} de
$1-p^+-p^->0$ (cette quantité n'est pas nécessairement strictement positive) ;
elle découle des boucles réflexives aux deux bornes de l'espace d'états -- en
$k=1$, une poussée supplémentaire vers $\Exc$ laisse l'état inchangé
($\max(1,k-1)=1$), et en $k=2N$, une poussée supplémentaire vers $\Inc$ laisse
l'état inchangé -- ce qui donne un état apériodique, et donc, par
irréductibilité, une chaîne entièrement apériodique. Unicité de $\pi$ par le
théorème ergodique ; forme géométrique par bilan détaillé
($\pi(k)p^+=\pi(k{+}1)p^- \Rightarrow \pi(k{+}1)/\pi(k)=\rho$ constant) ;
somme géométrique pour $\pi(\Inc)=\sum_{k>N}\pi(k)$. La convergence p.s. de la
moyenne temporelle est le théorème ergodique standard appliqué à l'indicatrice
$\mathbf 1[\cdot>N]$. Aucune hypothèse sur la forme de $P(X,Y)$ n'est utilisée
dans cette preuve : elle ne dépend que du rapport $\rho=p^+/p^-$, quelle que
soit la loi jointe sous-jacente.
\end{proof}

\begin{note}[Terminologie : convergence en distribution, pas absorption]
Le terme \emph{absorption} est réservé au cas unidirectionnel
(Corollaire~\ref{cor:unidirectionnel}, $p^-=0$ ou $p^+=0$). Dans le cas
ergodique ($p^+,p^->0$), l'automate ne se fixe \emph{jamais} définitivement :
il continue de se déplacer indéfiniment et converge seulement \emph{en
distribution} vers $\pi$. « Sans bruit » n'implique donc pas absorption : pour
une conjonction comme $Y=X_1\wedge X_2$, un automate peut recevoir des
poussées dans les deux sens même en l'absence de bruit sur $Y$, car les autres
combinaisons de $(X_1,X_2)$ jouent un rôle symétrique à celui du bruit.
\end{note}

\begin{corollary}[Cas unidirectionnel absorbant]\label{cor:unidirectionnel}
Trois cas distincts, à distinguer explicitement :
\begin{itemize}
  \item $p^-=0,\ p^+>0$ : chaîne de naissance pure (aucune poussée vers
  $\Exc$), absorbée p.s.\ en $2N$ en temps fini, $\rho\to\infty$. Deux temps
  d'arrêt distincts sont pertinents et \emph{ne doivent pas être confondus} :
  le temps d'\emph{absorption complète} en $2N$,
  $\mathbb E[T_{2N}]=(2N-s_0)/p^+$ (loi négative binomiale, cf.\ Lemme~3.1 de
  \textit{Démonstration de convergence}), et le temps, généralement plus
  court, de \emph{première entrée dans la région $\Inc$}
  (franchissement de $N$), pertinent pour le comportement du classifieur :
  $\mathbb E[T_{\Inc}]=(N+1-s_0)_+/p^+$ (où $(\cdot)_+=\max(\cdot,0)$, nul si
  $s_0>N$, l'automate étant déjà inclus) ;
  \item $p^+=0,\ p^->0$ : symétrique, absorption p.s.\ en $1$,
  $\mathbb E[T_1]=(s_0-1)/p^-$, première entrée dans $\Exc$ (franchissement de
  $N{+}1$) en $\mathbb E[T_{\Exc}]=(s_0-N)_+/p^-$ ;
  \item $p^+=p^-=0$ : l'automate ne reçoit \emph{aucun} signal ; chaque état
  est trivialement absorbant au sens où aucune transition ne se produit
  jamais (et non « non absorbant » -- il n'y a simplement pas de dynamique du
  tout). L'automate reste figé à son état initial $s_0$, cas exclu du
  Théorème~\ref{thm:multibit}.
\end{itemize}
Les deux premiers cas sont les cas particuliers $\rho=\infty$ et $\rho=0$ du
Théorème~\ref{thm:stationnaire}.
\end{corollary}

% ======================================================
\section{Lemme de non-contradiction}
% ======================================================

\begin{lemma}[Exclusion des dérives simultanément positives]\label{lem:noncontradiction}
Notons $\Delta = p^+-p^-$ et $\bar\Delta=\bar p^+-\bar p^-$ (les drifts nets de
$v$ et $\vbar$). Alors, pour toute loi $P(X,Y)$ :
\[
  \boxed{\Delta+\bar\Delta = -a\cdot P(Y{=}1) \;\le\; 0.}
\]
\end{lemma}

\begin{proof}
D'après \eqref{eq:pv}--\eqref{eq:pvbar} :
\begin{align*}
  \Delta &= S\,P(0,0) - S\,P(1,0) - a\,P(0,1), \\
  \bar\Delta &= S\,P(1,0) - S\,P(0,0) - a\,P(1,1).
\end{align*}
En sommant, les termes en $S\,P(0,0)$ et $S\,P(1,0)$ s'annulent deux à deux :
\[
  \Delta+\bar\Delta = -a\bigl[P(0,1)+P(1,1)\bigr] = -a\,P(Y{=}1) \le 0
\]
puisque $a>0$ et $P(Y{=}1)\ge0$.
\end{proof}

\begin{note}[Interprétation : exclusion asymptotique, pas instantanée]
$v$ et $\vbar$ ne peuvent donc jamais avoir \emph{simultanément} un drift net
strictement positif ($\Delta>0$ et $\bar\Delta>0$) : la règle empêche la
configuration « contradiction » ($x\wedge\neg x$, ligne 1 du tableau
sémantique) d'être la configuration \emph{asymptotiquement favorisée}. Elle
n'empêche en revanche pas les deux drifts d'être simultanément
\emph{négatifs} ($\Delta<0$ et $\bar\Delta<0$, régime neutre du
Théorème~\ref{thm:selection}) : le lemme exclut seulement la double
inclusion, pas la double exclusion. Il ne dit rien non plus sur l'état à un
instant $t$ fini : par le Théorème~\ref{thm:stationnaire}, chaque automate
garde une probabilité stationnaire résiduelle de se trouver « du mauvais
côté », décroissant exponentiellement en $N$ \emph{dès lors que $\rho\ne1$}
(hors régime frontière -- à $\rho=1$, cette probabilité reste exactement
$\tfrac12$, cf.\ Note « Régime frontière », Section~4) mais jamais strictement
nulle pour $N$ fini. La contradiction transitoire (les deux automates
simultanément en $\Inc$) n'est donc pas structurellement impossible à tout
instant ; elle est seulement exponentiellement rare pour $N$ grand hors
régime frontière. L'égalité $\Delta+\bar\Delta=0$ n'a lieu que si $P(Y{=}1)=0$
(aucun exemple positif) : cas dégénéré déjà identifié comme limite
(\textit{Démonstration de convergence}, Remarque sur $c=1$).
\end{note}

\begin{corollary}[Asymétrie stricte entre littéral et négation]\label{cor:asymetrie}
Si $\Delta>0$, alors $\bar\Delta = -a\,P(Y{=}1) - \Delta < 0$
strictement (et réciproquement, par symétrie du Lemme~\ref{lem:noncontradiction}).
Dès que $x$ est asymptotiquement favorisé, $\neg x$ est \emph{strictement}
défavorisé -- pas seulement « non favorisé ».
\end{corollary}

% ======================================================
\section{Critère de sélection sous distribution arbitraire}
% ======================================================

\begin{theorem}[Sélection, cas général -- régimes non dégénérés]\label{thm:selection}
Sous la règle à 4 cas avec $S,a>0$ fixés, supposons les observations
$(X_t,Y_t)_{t\ge1}$ i.i.d.\ de loi $P$, et $\Delta\ne0$, $\bar\Delta\ne0$
(hypothèse de non-dégénérescence, cf.\ Note ci-dessous pour le cas frontière).
Notons $\gamma=\max(\rho,1/\rho)>1$ et $\bar\gamma=\max(\bar\rho,1/\bar\rho)>1$
(de sorte que, par le Théorème~\ref{thm:stationnaire}, la probabilité que $v$
--resp.\ $\vbar$-- soit du mauvais côté est toujours $\le\gamma^{-N}$ --resp.\
$\bar\gamma^{-N}$--, quel que soit le signe de $\Delta$). Alors :
\begin{align}
  \Delta>0\ (\rho>1,\ \pi(\Inc)\to1) &\iff S\,P(0,0) \;>\; S\,P(1,0)+a\,P(0,1); \label{eq:crit_v}\\
  \bar\Delta>0\ (\bar\rho>1,\ \bar\pi(\Inc)\to1) &\iff S\,P(1,0) \;>\; S\,P(0,0)+a\,P(1,1). \label{eq:crit_vbar}
\end{align}
Par le Corollaire~\ref{cor:asymetrie}, \eqref{eq:crit_v} et \eqref{eq:crit_vbar}
sont mutuellement exclusifs, et sous l'hypothèse de non-dégénérescence les
trois régimes suivants sont exhaustifs (on énonce directement la convergence
en probabilité du couple $(v_t,\vbar_t)$, \emph{sans} invoquer une loi
stationnaire jointe -- la chaîne jointe $(v_t,\vbar_t)_t$ n'est pas
nécessairement irréductible, cf.\ Remarque ci-après) :
\begin{itemize}
  \item \textbf{Régime de sélection de $x$} : $\Delta>0$, donc
  (Corollaire~\ref{cor:asymetrie}) $\bar\Delta<0$ $\Rightarrow$
  $\lim_{t\to\infty}\mathbb P\bigl((v_t,\vbar_t)\ne(\Inc,\Exc)\bigr)
  \le\gamma^{-N}+\bar\gamma^{-N}\xrightarrow[N\to\infty]{}0$ :
  la structure $x$ est asymptotiquement préférée.
  \item \textbf{Régime de sélection de $\neg x$} : $\bar\Delta>0$, donc $\Delta<0$
  $\Rightarrow$
  $\lim_{t\to\infty}\mathbb P\bigl((v_t,\vbar_t)\ne(\Exc,\Inc)\bigr)
  \le\gamma^{-N}+\bar\gamma^{-N}\xrightarrow[N\to\infty]{}0$ :
  la structure $\neg x$ est asymptotiquement préférée.
  \item \textbf{Régime d'exclusion des deux littéraux} : $\Delta<0$ \emph{et}
  $\bar\Delta<0$ $\Rightarrow$
  $\lim_{t\to\infty}\mathbb P\bigl((v_t,\vbar_t)\ne(\Exc,\Exc)\bigr)
  \le\gamma^{-N}+\bar\gamma^{-N}\xrightarrow[N\to\infty]{}0$ :
  la structure tautologique est asymptotiquement préférée, $x$ est ignoré.
\end{itemize}
La borne $\gamma^{-N}+\bar\gamma^{-N}$ sur le $\lim_t$ n'est valable qu'une
fois la chaîne proche de sa loi stationnaire (ou si elle y est initialisée) --
elle ne dit rien sur $\mathbb P((v_t,\vbar_t)\ne\cdot)$ à un instant $t$ petit
et une initialisation arbitraire, d'où le $\lim_{t\to\infty}$ et non une
égalité pour tout $t$. Pour tout $N$ fini, l'automate est ergodique et fluctue
selon $\pi_N$ ; c'est la limite $N\to\infty$ (après celle en $t$) qui donne la
sélection presque certaine (cf.\ Note terminologique, Section~2). Sous une loi
$P(X,Y)$ arbitraire, la structure ainsi préférée n'est pas nécessairement la
clause génératrice ni la clause Bayes-optimale -- une variable peut être
sélectionnée parce qu'elle est corrélée à une variable pertinente sans
appartenir au concept cible.
\end{theorem}

\begin{proof}
\eqref{eq:crit_v} : $\Delta>0 \iff p^+>p^- \iff S\,P(0,0) > S\,P(1,0)+a\,P(0,1)$
directement depuis \eqref{eq:pv}. Identique pour \eqref{eq:crit_vbar}.
L'exclusion mutuelle et l'implication $\Delta>0\Rightarrow\bar\Delta<0$
découlent du Corollaire~\ref{cor:asymetrie}. Sous $\Delta,\bar\Delta\ne0$, les
quatre combinaisons de signes $(\mathrm{sgn}\,\Delta,\mathrm{sgn}\,\bar\Delta)$
se réduisent aux trois listées, la combinaison $(+,+)$ étant exclue par le
corollaire. Passage de la convergence marginale à la convergence jointe (par
exemple pour le régime IDENTITY, $\Delta>0$, donc $\rho>1$ et $\bar\rho<1$,
soit $\gamma=\rho$ et $\bar\gamma=1/\bar\rho$) : en prenant d'abord
$\lim_{t\to\infty}$ (le Théorème~\ref{thm:stationnaire} borne la loi
\emph{stationnaire}, atteinte seulement asymptotiquement en $t$ depuis une
initialisation arbitraire -- l'inégalité ci-dessous serait fausse pour un $t$
petit), puis par union bound,
\[
  \lim_{t\to\infty}\mathbb P\bigl((v_t,\vbar_t)\ne(\Inc,\Exc)\bigr)
  \le \lim_{t\to\infty}\mathbb P(v_t\ne\Inc) + \lim_{t\to\infty}\mathbb P(\vbar_t\ne\Exc)
  \le \gamma^{-N} + \bar\gamma^{-N} \xrightarrow[N\to\infty]{} 0,
\]
sans jamais supposer l'indépendance de $v_t$ et $\vbar_t$ (l'union bound ne
requiert que les bornes marginales du Théorème~\ref{thm:stationnaire},
valables même si $v_t$ et $\vbar_t$ sont corrélés). Les deux autres régimes se traitent de
façon identique. \emph{Attention} : $\bar\gamma^{-N}=\bar\rho^{N}$ ici (et non
$\bar\rho^{-N}$, puisque $\bar\rho<1$) -- c'est pour éviter précisément cette
confusion de signe que $\gamma=\max(\rho,1/\rho)$ est défini une fois pour
toutes, indépendamment du sens de la préférence.
\end{proof}

\begin{remark}[Non-irréductibilité possible de la chaîne jointe]
La chaîne $(v_t,\vbar_t)_t$ est bien une chaîne de Markov homogène (les deux
composantes sont mises à jour par la même observation $(X_t,Y_t)$ i.i.d.\ à
chaque pas), mais elle n'est pas nécessairement irréductible sur
$\{1,\ldots,2N\}^2$ tout entier : par exemple avec $S=1$ et $P(Y{=}1)=0$, une
poussée vers $\Inc$ sur $v$ s'accompagne toujours d'une poussée vers $\Exc$ sur
$\vbar$ (et réciproquement), de sorte que $v_t+\vbar_t$ est conservé tant
qu'aucun des deux automates n'atteint une borne -- la chaîne jointe reste alors
confinée à une sous-classe déterminée par l'initialisation. C'est précisément
pour cette raison que la preuve ci-dessus évite d'invoquer une loi stationnaire
jointe $\Pi_N$ et travaille directement avec $\lim_t \mathbb
P((v_t,\vbar_t)\ne\cdot)$, borné via l'union bound sur les deux marginales
seules -- un argument qui ne requiert aucune irréductibilité de la chaîne
jointe.
\end{remark}

\begin{note}[Régime frontière $\Delta=0$ ou $\bar\Delta=0$]
Si l'hypothèse de non-dégénérescence est levée, $\Delta=0$ donne $\rho=1$ et,
par le Théorème~\ref{thm:stationnaire}, $\pi(\Inc)=\tfrac12$ : l'automate
correspondant n'a stationnairement \emph{aucune} préférence -- ni convergence
vers $\Inc$, ni vers $\Exc$ -- contrairement à ce qu'impliquerait une lecture
naïve de « les deux inégalités sont fausses ». Ce cas frontière correspond à
une égalité exacte entre deux quantités continues de $P(X,Y)$ et est donc de
mesure nulle pour une loi jointe générique, mais il doit être exclu
explicitement (hypothèse $\Delta,\bar\Delta\ne0$) pour que les trois régimes du
Théorème~\ref{thm:selection} soient effectivement exhaustifs et mutuellement
exclusifs.
\end{note}

\begin{note}[Unification de deux résultats déjà établis]
Ce théorème unique remplace et généralise à la fois le résultat sur le
« littéral pertinent » (\textit{Convergence de la clause conditionnée},
Théorème~3.1, cas $p_1\ne p_0$ dans la notation $(c,p_1,p_0)$ de cette
référence, où elle écrit $a,b$) et celui sur la « neutralité » (\emph{ibid.},
Théorème~4.2, cas $p_1{=}p_0{=}p_y$ restreint) : les deux étaient des cas
particuliers de la même paire d'inégalités \eqref{eq:crit_v}--\eqref{eq:crit_vbar},
simplement exprimées dans la paramétrisation $(c,p_1,p_0)$ plutôt qu'en
$P(X,Y)$ général. Le seuil $S_{\text{bas}}(j)$ et le seuil $S_i^{\max}$ déjà publiés se
retrouvent en isolant $S$ dans \eqref{eq:crit_v} pour les deux régimes
$P(0,1)>0$ (seuil bas) ou distribution symétrique (seuil haut) --- calcul
inchangé, notation unifiée.
\end{note}

% ======================================================
\section{Extension multi-bit : borne de concentration pour une clause}
% ======================================================

\begin{theorem}[Concentration multi-bit]\label{thm:multibit}
Soit une clause à $n$ features, dont $K$ fixées par la condition ; les $d=n-K$
restantes sont gouvernées par $2d$ automates. Pour chaque automate $j$, en
supposant sa préférence stationnaire non dégénérée, notons
$r_j=p_j^+/p_j^-$, $\gamma_j=\max(r_j,1/r_j)>1$, et
$\gamma_{\min}=\min_{1\le j\le 2d}\gamma_j$. Notons $C^\star$ la structure
préférée sous $P$ (Théorème~\ref{thm:selection}) et $C_t$ la clause
effectivement calculée à l'instant $t$. Alors, sans hypothèse d'indépendance
entre automates :
\[
  \boxed{\lim_{t\to\infty}\mathbb P\bigl(C_t\ne C^\star\bigr)
  \;\le\; \sum_{j=1}^{2d} \frac1{1+\gamma_j^N}
  \;\le\; 2(n-K)\cdot\gamma_{\min}^{-N}.}
\]
\end{theorem}

\begin{proof}
Si $C_t\ne C^\star$, au moins un des $2(n-K)$ automates est du mauvais côté.
Le résultat suit de la borne stationnaire à un seul automate
(Théorème~\ref{thm:stationnaire}) et de l'union bound, qui ne requiert aucune
indépendance entre les automates (juste leurs bornes marginales).
\end{proof}

\begin{note}[Cas exclus et portée du résultat]
Le cas frontière $r_j=1$ et le cas dégénéré $p_j^+=p_j^-=0$ (aucun signal) sont
exclus par l'hypothèse de non-dégénérescence, en cohérence avec
$\Delta,\bar\Delta\ne0$ du Théorème~\ref{thm:selection} ; un automate absorbant
(Corollaire~\ref{cor:unidirectionnel}) est couvert par la convention
$\gamma_j=\infty$. La borne utilise nécessairement le \emph{pire} automate
($\gamma_{\min}$), pas le meilleur. Enfin, $C_t\ne C^\star$ est un
\emph{mismatch structurel} (deux structures différentes peuvent calculer la
même fonction logique par coïncidence), pas nécessairement une erreur de
classification ; et la borne porte sur la clause interne conditionnelle, la
condition elle-même étant fixée par construction (sa couverture fait l'objet d'un
résultat séparé, Section~6).
\end{note}

\begin{note}[Lien avec le conditionnement $K$]
La loi $P(X,Y)$ de la Section~1 est ici la loi conditionnelle sur le
sous-ensemble applicable $D_F$ (les exemples satisfaisant la condition sur les $K$
features gelées). Sans conditionnement ($K=0$, donc $d=n$), le théorème
s'applique tel quel à la clause entière -- mais rien ne garantit alors que
chaque littéral individuel tombe dans le régime « pertinent » plutôt que
« neutre » pour une cible à $m\ge2$ littéraux : c'est exactement la limite déjà
documentée (échec de \texttt{hybrid\_coord\_4case.cpp}, test $K{=}0$ sur Iris à
22,7\,\%, sous le hasard).
\end{note}

% ======================================================
\section{Couverture des conditions aléatoires : garantie d'ensemble}
% ======================================================

Il faut séparer nettement deux sources d'erreur bien distinctes : (1) la
probabilité de \emph{tirer une bonne condition} (propriété fixe, déterministe
une fois $F,v$ choisis) et (2) sachant la condition bonne, la probabilité que
la clause associée \emph{apprenne effectivement} la bonne structure
(propriété probabiliste, dépendant de $N$ et $t$ -- déjà bornée par le
Théorème~\ref{thm:multibit}). Cette section formalise la première, puis les
combine, reliant $M$, $K$, $n$ et $N$ au ratio $M/(2^K\binom{n}{K})$ déjà
observé empiriquement (Section~\ref{sec:validation}).

Il existe $C_K=2^K\binom{n}{K}$ conditions possibles : le choix de $K$
features parmi $n$, et une valeur binaire pour chacune.

\begin{definition}[Condition utile]\label{def:condition-utile}
Fixons un concept cible (un sous-ensemble de littéraux pertinents parmi les
$n$). Une condition $(F,v)$, $F\subseteq\{1,\ldots,n\}$, $|F|=K$, $v\in\{0,1\}^K$,
est dite \emph{utile} pour ce concept si, sous la loi conditionnelle
$P(X,Y\mid X_F=v)$, la règle sélectionne correctement (au sens du
Théorème~\ref{thm:selection}) les littéraux cibles parmi les $n-K$ features
non fixées. Notons $G^\star\ge1$ le nombre de conditions utiles parmi les
$C_K$ possibles (hypothèse d'identifiabilité minimale : au moins une
condition révèle le concept), et
\[
  q = \frac{G^\star}{C_K} = \frac{G^\star}{2^K\binom{n}{K}}
\]
la probabilité qu'une condition tirée uniformément soit utile.
\end{definition}

\begin{theorem}[Couverture des conditions utiles]\label{thm:coverage}
Si les $M$ clauses tirent leur condition indépendamment et uniformément
(garanti par le code : \texttt{shuffle(pool.begin(), pool.end(), rng)} donne
un sous-ensemble de $K$ features uniforme, et
\texttt{vals.push\_back(valDist(rng))} tire chaque valeur par un
pile-ou-face indépendant) :
\[
  \boxed{\mathbb P(\text{au moins une condition utile est tirée})
  = 1-(1-q)^M \;\ge\; 1-e^{-Mq}.}
\]
\end{theorem}

\begin{proof}
Chaque clause tire une condition utile avec probabilité $q$, indépendamment
des autres. La probabilité qu'aucune des $M$ clauses n'en tire une vaut
$(1-q)^M$ ; le complémentaire donne l'égalité, et $1-q\le e^{-q}$ donne la
forme exponentielle.
\end{proof}

\begin{corollary}[Garantie de présence d'une clause correcte]\label{cor:end-to-end}
Supposons que, pour toute condition utile, la probabilité (asymptotique en
$t$) que la clause apprise soit structurellement incorrecte soit bornée par
$\varepsilon_N=2(n-K)\gamma_{\min}^{-N}$ (Théorème~\ref{thm:multibit}). Alors,
par union bound sur les deux complémentaires (« aucune condition utile
tirée », Théorème~\ref{thm:coverage} ; et, pour une clause ayant reçu une
condition utile, « elle n'apprend pas la bonne structure », Théorème~\ref{thm:multibit}) :
\[
  \mathbb P(\text{l'ensemble contient au moins une clause correcte})
  \;\ge\; 1-e^{-Mq}-\varepsilon_N = 1-e^{-Mq}-2(n-K)\gamma_{\min}^{-N}.
\]
\end{corollary}

\begin{note}[Deux sources d'erreur distinctes, et ce que ça ne prouve pas]
Point important : ne pas appeler une clause « utile » avant son entraînement
-- \emph{la condition est utile} (propriété fixe de $P$ et du concept cible,
indépendante de $N$ et $t$), \emph{puis} la clause peut être apprise
correctement ou non (propriété probabiliste du Théorème~\ref{thm:multibit}).
Ce corollaire garantit seulement l'existence d'\emph{au moins une} clause
correcte dans l'ensemble, pas que le vote final (pondéré par $\alpha_m$,
Section~6) la retient effectivement ni qu'elle domine les autres clauses au
moment du classement. Combiner cette garantie de couverture avec une garantie
sur le poids relatif $\alpha_m$ des clauses correctes reste un travail futur.
\end{note}

% ======================================================
\section{Analyse par round sous les distributions induites par le boosting}
% ======================================================

\begin{theorem}[Caractérisation conditionnelle d'un round de boosting -- distribution d'entraînement effective $Q_m$]\label{thm:boosting}
Soit, au round $m$, $S_m$ le sous-ensemble applicable de la clause (les
exemples satisfaisant son conditionnement) et $w_i$ les poids AdaBoost
courants sur $S_m$. Notons $W_m^+=\sum_{i\in S_m:\,y_i=1}w_i$ et
$W_m^-=\sum_{i\in S_m:\,y_i=0}w_i$. Le code d'entraînement
(\texttt{modele\_xor.cpp}, \texttt{modele\_multiclasse.cpp}) ne tire
\emph{pas} directement selon la distribution pondérée
$D_m(i)=w_i/\sum_{j\in S_m}w_j$ : il choisit d'abord la classe positive ou
négative avec probabilité $\tfrac12$ chacune \emph{si les deux sont non vides},
sinon il échantillonne exclusivement dans l'unique classe présente ; si $S_m$
est vide, la clause n'est pas entraînée du tout ($\alpha_m=0$,
\texttt{if (subsetIdx.empty()) \{ alphas.push\_back(0.0); continue; \}}). La
distribution réellement échantillonnée pour mettre à jour les automates est
donc
\[
  Q_m(i) = \begin{cases}
    \tfrac12\cdot\dfrac{w_i}{W_m^+}, & i\in S_m,\ y_i=1,\ W_m^+>0,\ W_m^->0,\\[2mm]
    \tfrac12\cdot\dfrac{w_i}{W_m^-}, & i\in S_m,\ y_i=0,\ W_m^+>0,\ W_m^->0,\\[2mm]
    \dfrac{w_i}{W_m^+}, & i\in S_m,\ y_i=1,\ W_m^->0 \text{ faux (sous-ensemble pur positif)},\\[2mm]
    \dfrac{w_i}{W_m^-}, & i\in S_m,\ y_i=0,\ W_m^+>0 \text{ faux (sous-ensemble pur négatif)},\\[2mm]
    0, & i\notin S_m,
  \end{cases}
\]
non définie (clause ignorée, $\alpha_m=0$) si $W_m^+=W_m^-=0$ (i.e.\ $S_m=\emptyset$).
Conditionnellement à l'historique du boosting et au tirage aléatoire du stump
(sous-ensemble conditionnant) du round $m$, et sur l'évènement $S_m\ne\emptyset$,
les Théorèmes~\ref{thm:stationnaire}, \ref{thm:selection}, \ref{thm:multibit}
s'appliquent à la clause $m$, \emph{chacun sous ses hypothèses respectives},
avec $P(X,Y)$ remplacée par $Q_m(X,Y)$ -- \emph{et non} par $D_m$. En
particulier, sur un sous-ensemble pur ($W_m^+=0$ ou $W_m^-=0$) ou
proche de la frontière ($r_j\approx1$ pour certains automates $j$), certains
littéraux de la clause $m$ peuvent relever du Corollaire~\ref{cor:unidirectionnel}
(absorbant) plutôt que du régime ergodique du Théorème~\ref{thm:stationnaire},
ou se trouver dans le régime frontière exclu du Théorème~\ref{thm:selection} ;
le Théorème~\ref{thm:multibit} ne s'applique alors qu'aux littéraux non
dégénérés de la clause $m$.
\end{theorem}

\begin{proof}
Les probabilités de transition \eqref{eq:pv}--\eqref{eq:pvbar} sont des
fonctions linéaires de la loi utilisée pour générer les tirages effectivement
présentés à l'automate ; substituer $Q_m$ à $P$ ne change aucune étape des
preuves précédentes, qui ne supposent jamais une forme particulière de $P$.
$Q_m$, et non $D_m$, est la loi qui génère les $(x,y)$ réellement vus par
\texttt{updateMasked} pendant le round $m$ : la ligne
\texttt{useNeg = subNeg.empty() ? false : (subPos.empty() ? true : !flip(0.5))}
donne exactement les quatre cas ci-dessus (rééquilibré si les deux classes sont
présentes, sinon échantillonnage pur dans l'unique classe présente). $D_m$
n'intervient nulle part dans cette mise à jour, uniquement dans le calcul de
l'erreur pondérée (\texttt{err += w[i]}, puis \texttt{err /= wSum}) et de
$\alpha_m$.
\end{proof}

\begin{note}[$Q_m$ pour l'apprentissage, $D_m$ pour l'évaluation]
Le rééquilibrage 50/50 (tirer selon $Q_m$ plutôt que $D_m$) est un choix de
conception délibéré : il évite qu'une classe majoritaire domine l'entraînement
de la clause. Mais l'erreur pondérée
$\mathrm{err}=\bigl(\sum_i w_i\mathbf 1[\text{faux}]\bigr)/\sum_i w_i$ et le
coefficient $\alpha_m=\tfrac12\ln\frac{1-\mathrm{err}}{\mathrm{err}}$ sont eux
calculés sous $D_m$, la vraie distribution AdaBoost non rééquilibrée. Le
Théorème~\ref{thm:boosting} caractérise donc la clause $m$ elle-même sous
$Q_m$ ; la qualité du weak learner (son erreur, son poids $\alpha_m$ dans le
vote final) est mesurée sous $D_m$. Ce n'est pas une incohérence du code, mais
une distinction nécessaire entre « ce qui entraîne la clause » et « ce qui
évalue sa qualité » -- absente d'une version antérieure de ce document, qui
remplaçait $P$ par $D_m$ partout, y compris pour caractériser l'entraînement
lui-même.
\end{note}

\begin{note}[Ce que ce théorème ne prouve pas]
Il caractérise chaque round \emph{isolément}, conditionnellement à $Q_m$ ---
il ne prouve pas la convergence de l'ensemble $F_c(x)=\sum_m
\alpha_{c,m}h_{c,m}(x)$ vers un classifieur de Bayes global, ce que
l'AdaBoost classique établit sous des hypothèses différentes (voir Related
Work). Nous ne revendiquons pas ce résultat plus fort.
\end{note}

\begin{note}[Weak learner abstentionniste]
Le weak learner $h_m$ associé à la clause $m$ n'est pas un classifieur binaire
standard : il s'abstient (vote nul) hors de $S_m$. En notation $\{-1,0,+1\}$,
$h_m(x)=0$ pour $x\notin S_m$. L'erreur pondérée
$\mathrm{err}_m=\bigl(\sum_{i\in S_m}w_i\mathbf1[h_m(x_i)\ne y_i]\bigr)/\sum_{i\in S_m}w_i$
et $\alpha_m=\tfrac12\ln\frac{1-\mathrm{err}_m}{\mathrm{err}_m}$ sont donc les
formules standard d'AdaBoost \emph{à classifieurs abstentionnistes}
(restreintes à $S_m$), pas celles d'AdaBoost classique sur l'ensemble complet
des exemples.
\end{note}

\begin{note}[Écart entre la structure caractérisée et le weak learner conservé]
Les Théorèmes~\ref{thm:selection} et \ref{thm:multibit} caractérisent la
structure $C^\star$ vers laquelle les automates de la clause $m$
convergent -- y compris le régime « exclusion des deux littéraux » pour
\emph{chacun} des $d-K$ littéraux libres, qui correspond à une clause
\emph{structurellement vide} si tous les littéraux libres y tombent
simultanément. Le code ne traite cependant pas une clause vide comme le
prescrit la convention usuelle des Tsetlin Machines (sortie $1$ à
l'entraînement, $0$ au test, cf.\ Granmo 2018, Eq.~6) : il la détecte après
coup (\texttt{isEmpty()}) et l'écarte purement et simplement du vote final
($\alpha_m=0$, clause ignorée). Les Théorèmes~\ref{thm:stationnaire}--\ref{thm:multibit}
décrivent donc la structure vers laquelle \emph{chaque} clause converge, mais
pas l'étape de sélection supplémentaire (garder/écarter) qui détermine
l'ensemble effectif de weak learners utilisés par le classifieur final ; cette
étape n'est pas couverte par les théorèmes ci-dessus et resterait à formaliser
séparément (par exemple en conditionnant explicitement sur l'évènement
« clause non vide »).
\end{note}

% ======================================================
\section{Validation empirique --- déjà réalisée}\label{sec:validation}
% ======================================================

La courbe $N\mapsto\pi(\Inc)=\rho^N/(1+\rho^N)$ (Théorème~\ref{thm:stationnaire})
et le seuil critique $\alpha^\star(S)$ issu du Théorème~\ref{thm:selection} ont
été confrontés à des données réelles sur \textbf{quatre} datasets
(\textit{Rôle de $S$ et $\alpha$ : théorie et validation} -- ce document utilise
encore la notation $\alpha$ pour ce même paramètre, ici renommé $a$ pour éviter
la collision avec $\alpha_m$) :
\begin{itemize}
  \item XOR : 42 configurations, accord exact sur la position du seuil (mur
  d'effondrement prédit à $\alpha^\star(S)=0{,}519S$, confirmé configuration par
  configuration) ;
  \item Iris, Digits, MNIST : confirment que la sensibilité observée dépend du
  ratio $M/(2^K\binom{n}{K})$ (nombre de clauses / nombre total de
  conditions possibles -- le choix des $K$ features \emph{et} d'une valeur
  binaire pour chacune, cf.\ \texttt{vals.push\_back(valDist(rng))} dans
  \texttt{modele\_multiclasse.cpp} : le nombre de combinaisons de features
  seul, $\binom{n}{K}$, sous-compte d'un facteur $2^K$), avec un test direct
  ($M{=}500\to16\,000$ sur MNIST) confirmant que l'écart se réduit quand ce
  ratio augmente.
\end{itemize}
Le classifieur multiclasses boosté (Théorème~\ref{thm:boosting}) a par ailleurs
déjà été comparé à SVM, Régression Logistique, Random Forest et MLP sur Iris,
Digits et MNIST, avec des résultats égalant ou dépassant ces baselines sur
Iris et Digits (voir document de suivi des expériences).

\end{document}
